\section{3GPP-Codebook und Beam-Codebook: zwei verschiedene $W$}
\label{sec:zwei-matrizen}

Es gibt zwei verschiedene \glqq$W$\grqq{} im Code. Die 3GPP-Matrix ist \code{make\_type1\_sp\_mode1()}.
\code{generate\_beam\_weights\_codebook()} ist \emph{nicht} die Umsetzung des 3GPP-Codebooks.

\subsection{Das 3GPP-Codebook}

\file{lib/ran/precoding/type1\_sp/precoding\_codebooks.cpp}, \code{make\_type1\_sp\_mode1()}, erzeugt aus dem PMI
die Matrix $W$ mit Dimension CSI-RS-Ports $\times$ Layer, genau wie in TS\,38.214 Tabelle 5.2.2.2.1-5 bis -8.
Pro Layer ein Vektor $[v_{l,m};\,\varphi_n v_{l,m}]$, normiert mit $1/\sqrt{\upsilon P}$.

\subsection{Das Beam-Codebook}

\file{lib/ran/beamforming/beam\_weights\_codebook\_generator.cpp} steht in keiner 3GPP-Spezifikation. Es ist eine
implementierungsspezifische Beam-Tabelle: alle DFT-Richtungen, die das Antennen-Array bilden kann, plus die
Einzelport-Beams. Das Konzept stammt aus O-RAN, wo die Radio Unit eine feste Tabelle von Beam-IDs kennt.
Zeilen sind Beams, Spalten physische Antennen.

\begin{lstlisting}[language=C++,title={lib/ran/beamforming/beam\_weights\_codebook\_generator.cpp}]
beam_weights_codebook ocudu::generate_beam_weights_codebook(antenna_topology topology)
{
  // ... Topologie-Parameter: nof_panels, nof_elements_dim1/2, nof_polarizations, nof_beams_dim1/2 ...
  beam_weights_codebook beam_weights = beam_weights_codebook(nof_total_beams, nof_total_antennas);

  // Direkte Port-Beams: Einheitsmatrix in den ersten nof_total_antennas Zeilen.
  for (unsigned i_antenna = 0; i_antenna != nof_total_antennas; ++i_antenna) {
    beam_weights.set_coefficient(1.0, get_beam_id(topology, i_antenna), i_antenna);
  }

  // Only direct beam to port mapping if the panels are 1x1 elements.
  if (!has_beam_grid(topology)) {
    return beam_weights;
  }

  float amplitude = std::sqrt(1.0F / static_cast<float>(nof_elements_dim1 * nof_elements_dim2));

  for (unsigned i_panel = 0; i_panel != nof_panels; ++i_panel) {
    for (unsigned i_pol = 0; i_pol != nof_polarizations; ++i_pol) {
      for (unsigned i_beam_dim1 = 0; i_beam_dim1 != nof_beams_dim1; ++i_beam_dim1) {
        for (unsigned i_beam_dim2 = 0; i_beam_dim2 != nof_beams_dim2; ++i_beam_dim2) {
          beam_identifier beam_id = get_beam_id(topology, i_panel, i_pol, i_beam_dim1, i_beam_dim2);
          unsigned i_port = nof_elements_dim1 * nof_elements_dim2 * (nof_polarizations * i_panel + i_pol);

          for (unsigned j = 0; j != nof_elements_dim1; ++j) {
            cf_t v_l = std::polar(amplitude, TWOPI * j * i_beam_dim1 / nof_beams_dim1);
            for (unsigned k = 0; k != nof_elements_dim2; ++k) {
              cf_t u_m   = std::polar(1.0F, TWOPI * k * i_beam_dim2 / nof_beams_dim2);
              beam_weights.set_coefficient(v_l * u_m, beam_id, i_port++);
            }
          }
        }
      }
    }
  }
  return beam_weights;
}
\end{lstlisting}

Ausgabe für \code{single\_panel\_two\_one} (20 Beams $\times$ 4 Antennen, gekürzt):
\begin{lstlisting}
  beam  0: [+1.00+0.00j  +0.00+0.00j  +0.00+0.00j  +0.00+0.00j]
  beam  3: [+0.00+0.00j  +0.00+0.00j  +0.00+0.00j  +1.00+0.00j]
  beam  4: [+0.71+0.00j  +0.71+0.00j  +0.00+0.00j  +0.00+0.00j]   l=0, pol 0
  beam  5: [+0.71+0.00j  +0.50+0.50j  +0.00+0.00j  +0.00+0.00j]   l=1, pol 0
  beam 12: [+0.00+0.00j  +0.00+0.00j  +0.71+0.00j  +0.71+0.00j]   l=0, pol 1
  beam 13: [+0.00+0.00j  +0.00+0.00j  +0.71+0.00j  +0.50+0.50j]   l=1, pol 1
\end{lstlisting}

\subsection{Die Zerlegung}

Die Steuervektoren $v_{l,m}$ des Type-I-Codebooks sind DFT-Vektoren mit Phase $2\pi l i/(N_1 O_1)$. Das
Beam-Codebook enthält dieselben Vektoren. Deshalb lässt sich jede 3GPP-Matrix zerlegen:
\begin{equation}
  W_{\text{3GPP}}\;(\text{Antennen}\times\text{Layer})
  = W_{\text{Beam}}^{\mathsf T}\;(\text{Antennen}\times\text{Beams})\cdot
    W_{\text{MIMO}}\;(\text{Beams}\times\text{Layer}).
\end{equation}
$W_{\text{MIMO}}$ ist winzig: pro Layer zwei Einträge, einer pro Polarisation mit $\pm\varphi$. Es entsteht in
\code{calculate\_mimo\_matrix()}:

\begin{lstlisting}[language=C++,title={lib/ran/precoding/type1\_sp/precoding\_codebooks.cpp, calculate\_mimo\_matrix()}]
precoding_beam_list beams = get_beams_from_pmi(pmi, nof_layers);   // Beam-IDs aus i_1_1, i_1_2, i_1_3
unsigned nof_beams = beams.size() / 2;                               // ohne Polarisationsdimension

unsigned i_2 = pmi.i_2 % pmi_ranges.i_2;
cf_t     phi = std::polar(1.0F, (TWOPI / 4.0F) * static_cast<float>(i_2));

precoding_weight_matrix weights(nof_layers, beams.size());

// Normalize precoding by the number of layers and ports, apply the reciprocal of the beamforming scaling.
float scaling = std::sqrt(static_cast<float>(info.n1 * info.n2)) / std::sqrt(static_cast<float>(nof_layers * nof_ports));

for (unsigned i_layer = 0; i_layer != nof_layers; ++i_layer) {
  unsigned i_beam    = i_layer % nof_beams;
  // First half of the layers uses +phi, second half uses -phi, as per TS38.214 Section 5.2.2.2.1.
  cf_t     layer_pol = (i_layer < (nof_layers + 1) / 2) ? phi : -phi;

  weights.set_coefficient(scaling,             i_layer, 2 * i_beam);      // Polarisation 0
  weights.set_coefficient(scaling * layer_pol, i_layer, 2 * i_beam + 1);  // Polarisation 1
}
return {weights, beams};
\end{lstlisting}

Der Precoding-Code kennt das Beam-Codebook nicht, nur die Beam-\emph{Nummern} über \code{get\_beam\_id()}.
Beide Seiten sind ausschließlich über diese Nummerierungskonvention gekoppelt. Dass sie zusammenpassen, sichert der
Test \file{precoding\_matrix\_factorization\_test}.

\paragraph{Zahlenbeispiel} \code{single\_panel\_two\_one}, 2 Layer, $i_{1,1}=1$, $i_{1,3}=0$, $i_2=1$:
beide Layer auf Beam $l=1$, Beam-IDs 5 (Pol.~0) und 13 (Pol.~1), $\varphi=j$, Skalierung
$\sqrt{2}/\sqrt{2\cdot4}=0{,}5$.
\begin{equation*}
  W_{\text{MIMO}} = \begin{bmatrix} 0{,}5 & +0{,}5j \\ 0{,}5 & -0{,}5j \end{bmatrix}
  \quad(\text{Layer}\times\text{Beams }5,13)
\end{equation*}
Layer 0 auf den Antennen: $0{,}5\cdot\text{beam}_5 + 0{,}5j\cdot\text{beam}_{13}
= [\,0{,}354,\; 0{,}25+0{,}25j,\; 0{,}354j,\; -0{,}25+0{,}25j\,]$, identisch mit der Spalte aus
\code{make\_type1\_sp\_mode1()}.

\subsection{Wo die Multiplikation mit den IQ-Samples passiert}

\begin{enumerate}
  \item \textbf{Upper PHY, Precoding.} Der Resource-Grid-Mapper (\file{lib/phy/support/resource\_grid\_mapper\_impl.cpp:127})
    ruft pro PRG \code{precoder.apply\_precoding()}. Mit Beam-Raster ist das $W_{\text{MIMO}}$, ohne die volle
    3GPP-Matrix. Die Rechnung steht in \file{lib/phy/generic\_functions/precoding/channel\_precoder\_generic.cpp}:
\begin{lstlisting}[language=C++]
for (unsigned i_re = 0; i_re != nof_re; ++i_re) {
  // Set the port RE to the contribution of the first layer.
  cf_t sum = layer_re_view_list[0][i_re] * port_weights[0];
  for (unsigned i_layer = 1; i_layer != nof_layers; ++i_layer) {
    // Accumulate the contributions of all other layers.
    sum += layer_re_view_list[i_layer][i_re] * port_weights[i_layer];
  }
  port_re[i_re] = sum;
}
\end{lstlisting}
  \item \textbf{Lower PHY, Beamforming.} Beim OFDM-Modulieren holt
    \file{lib/phy/lower/processors/downlink/pdxch/pdxch\_baseband\_modulator.h:163} für jede physische Antenne die
    Spalte des Beam-Codebooks und der Modulator bildet die gewichtete Summe über die Beam-Ports des Grids.
    Nullgewichte werden übersprungen (\code{std::isnormal}-Prüfung in
    \file{lib/phy/lower/modulation/ofdm\_modulator\_impl.cpp:86}).
\end{enumerate}

\subsection{Wann \code{generate\_beam\_weights\_codebook()} aufgerufen wird}

Genau einmal pro Zelle beim Start, unabhängig vom Traffic, in der Initialisierungsliste von
\code{pdxch\_processor\_impl} (\file{lib/phy/lower/processors/downlink/pdxch/pdxch\_processor\_impl.h:52}) sowie für
Open Fronthaul in \file{lib/ofh/transmitter/ofh\_transmitter\_factories.cpp:68}.

\begin{lstlisting}[language=C++,title={lib/phy/lower/processors/downlink/pdxch/pdxch\_processor\_impl.h}]
pdxch_processor_impl(std::unique_ptr<ofdm_symbol_modulator> modulator,
                     task_executor&                         executor,
                     const configuration&                   config) :
  logger(ocudulog::fetch_basic_logger("PHY")),
  beamforming_codebook(generate_beam_weights_codebook(config.tx_ant_topology)),
  ...
\end{lstlisting}

Ohne \code{tx\_ant\_topology} in der Konfiguration wird die Topologie aus \code{nof\_antennas\_dl} abgeleitet:
\code{one\_port} bis \code{eight\_ports}. Keine davon hat ein Beam-Raster, das Codebook ist dann die Einheitsmatrix
und wirkungslos.

\begin{table}[h]\centering
\begin{tabular}{lll}
\toprule
Setup & Beamforming geschieht in & Beam-Codebook \\
\midrule
SDR, \code{four\_ports} & Upper PHY, volle 3GPP-Matrix & Einheitsmatrix, wirkungslos \\
SDR, \code{single\_panel\_two\_one} & Upper PHY $W_{\text{MIMO}}$, Lower PHY Tabelle & nötig \\
OFH-RU mit Beamforming & Upper PHY $W_{\text{MIMO}}$, RU Tabelle & nötig, wird an RU gesendet \\
\bottomrule
\end{tabular}
\end{table}

\subsection{Wofür das Beam-Codebook trotzdem existiert}

\begin{itemize}
  \item \textbf{Open Fronthaul mit Beamforming in der Radio Unit.} Die DU sendet pro Beam-Port und nennt im
    C-Plane die Beam-ID. \file{lib/ofh/support/beamforming\_weights\_repository.cpp} komprimiert jede Zeile der
    Tabelle einmal und schickt sie an die RU.
  \item \textbf{Beams ohne PMI.} SSB, PDCCH, SIB und CSI-RS haben kein UE, das einen PMI meldet. Die
    Zellkonfiguration zeigt mit \code{ref\_beams} direkt auf Zeilen der Tabelle (z.\,B. SSB-Sweeping).
  \item \textbf{Ein Grid-Format für beide Backends.} Der Upper PHY schreibt immer dasselbe Resource Grid mit
    Beam-Ports, egal ob dahinter die eigene Lower PHY oder eine OFH-RU hängt.
\end{itemize}

\subsection{Der UHD-SDR-Fall}

Mit einem UHD-SDR und Standardtopologie passiert das komplette Beamforming im Upper PHY mit der 3GPP-Matrix:
\begin{enumerate}
  \item Codebook beim Start: \code{make\_type1\_sp\_mode1()} für jeden PMI, FAPI-Repository speichert.
  \item Precoding pro PDSCH: \code{apply\_precoding\_port()}, pro Port eine Spalte von $W$, pro Resource Element
    eine Summe über die Layer. Das ist $W\cdot x$ auf Subträgerebene, vor der IFFT.
  \item OFDM-Modulation: pro Port IFFT und Cyclic Prefix. Beam-Codebook ist Einheitsmatrix.
  \item UHD: die Antennen-Zeitsignale gehen als TX-Kanäle an das Radio. Die relative Phase zwischen den Kanälen ist
    die Keule; die Hardware muss phasenkohärent senden (gemeinsamer LO, z.\,B. N310/X310).
\end{enumerate}
Konfiguration: \code{nof\_antennas\_dl: 4}, \code{tx\_ant\_topology} nicht setzen. Das Codebook ist dann
\code{two\_one}, 2 kreuzpolarisierte Elemente, 8 Richtungen über $i_{1,1}$. Für 16 Richtungen mit 4 Elementen pro
Polarisation sind 8 TX-Kanäle nötig (\code{four\_one}).

\subsection{Was OCUDU beim Beamforming nicht macht}

\begin{itemize}
  \item Keine Keule auf einen beliebigen Winkel in Grad, nur das feste DFT-Raster.
  \item Keine Nebenkeulenunterdrückung durch Amplituden-Tapering; alle Koeffizienten haben denselben Betrag.
  \item Kein Downlink-Beamforming aus SRS-Messungen; die SRS-Kanalmatrix dient nur der Uplink-TPMI-Auswahl.
  \item Keine Nullstellen auf andere UEs, kein Zero-Forcing für Multi-User-MIMO.
\end{itemize}
Die Richtung bestimmt das UE über den PMI-Report; OCUDU wendet sie an. Im Testmodus setzt man sie mit $i_{1,1}$ fest.
